% ============================================================
%  Introducing Python, chapter 27: Performance
% ============================================================
\chapter{Performance}

\tikzset{
  w27n/.style={draw=#1,fill=#1!10,rounded corners=2pt,inner sep=4pt,
    minimum height=7mm,font=\sffamily\small,align=center},
  w27n/.default=cblue
}

\textbf{What this chapter gives you.} Python is fast enough until it is not, and people guess wrong
about where the time goes. This chapter teaches you to \emph{measure} first (timers and a
profiler), then fix the real hotspot with a better algorithm, a cache, concurrency, or a faster
tool. In AI work the slow parts are usually waiting on model APIs, big loops over data, and
memory use.

%MD
%```mermaid
%flowchart LR
%    M["1. Measure (timer / profiler)"] --> F["2. Find the real hotspot"] --> C["3. Change ONE thing"] --> R["4. Re-measure: faster? keep it"]
%    R -.->|"repeat while it is still too slow"| M
%```
%ENDMD
\begin{center}
\begin{tikzpicture}[font=\sffamily\small,>=Stealth]
  \node[w27n=cblue]   (m) at (0,0)    {1. Measure\\{\scriptsize timer / profiler}};
  \node[w27n=cred]    (f) at (4.0,0)  {2. Find\\{\scriptsize the real hotspot}};
  \node[w27n=cgold]   (c) at (8.0,0)  {3. Change\\{\scriptsize ONE thing}};
  \node[w27n=cgreen]  (r) at (12.0,0) {4. Re-measure\\{\scriptsize faster? keep it}};
  \foreach \a/\b in {m/f,f/c,c/r} \draw[->,thick] (\a) -- (\b);
  \draw[->,thick,dashed,cgray] (r.south) -- ++(0,-0.7) -| (m.south);
  \node[font=\scriptsize,text=cgray,fill=white] at (6.0,-0.7) {repeat while it is still too slow; stop when it is fast enough};
\end{tikzpicture}
\end{center}

% ------------------------------------------------------------
\section{Computer Hardware}

A CPU core has registers (hold the numbers in use), an arithmetic-logic unit (does the
operations), a control unit (fetches instructions and data) and a clock. Around it sit caches, RAM,
an SSD and maybe a hard drive. Going down the list, storage gets bigger, cheaper and \textbf{much
slower}. The first two kinds lose their data when power is lost; the last two keep it.

%MD
%```text
%Time for one access (rough orders of magnitude, illustrative, log scale)
%Register          | #                    under 1 ns
%L1 cache          | #                    about 1 ns
%L2 cache          | ##                   few ns
%L3 cache          | ###                  about 10 ns
%RAM               | #####                about 100 ns
%SSD               | ################     about 0.1 ms
%Hard disk         | ###################  several ms
%Network API call  | #####################  100 ms and up
%                    1 ns     1 us      1 ms      1 s
%```
%ENDMD
\begin{center}
\begin{tikzpicture}
\begin{axis}[xbar,xmode=log,xmin=100,xmax=1e14,width=13cm,height=7cm,
  symbolic y coords={Register,L1 cache,L2 cache,L3 cache,RAM,SSD,Hard disk,Network API call},
  ytick=data,y dir=reverse,xlabel={time for one access (log scale)},
  bar width=8pt,enlarge y limits=0.1,nodes near coords,point meta=explicit symbolic,
  every node near coord/.append style={font=\scriptsize,anchor=west},
  xtick={1e3,1e6,1e9,1e12},xticklabels={1 ns,1 $\mu$s,1 ms,1 s},
  title={\scriptsize rough orders of magnitude, illustrative: real numbers vary by machine},title style={font=\sffamily}]
\addplot[fill=cblue!60,draw=cblue] coordinates {
  (300,Register) [under 1 ns] (1000,L1 cache) [about 1 ns] (4000,L2 cache) [few ns]
  (12000,L3 cache) [about 10 ns] (80000,RAM) [about 100 ns] (1e8,SSD) [about 0.1 ms]
  (5e9,Hard disk) [several ms] (1e11,Network API call) [100 ms and up]};
\end{axis}
\end{tikzpicture}
\end{center}

The point of the chart: each step down costs 10 to 1000 times more. Waiting on a network call is
millions of times slower than adding two numbers. The operating system moves data between rows for
you (for example virtual memory pages to disk, which is why a computer that runs out of RAM
suddenly crawls). At the software level you can help in three ways:

\begin{center}
\begin{tabularx}{\linewidth}{@{}l X X@{}}
\toprule
\textbf{Lever} & \textbf{Idea} & \textbf{Python example} \\
\midrule
Data layout & contiguous data fits caches; scattered objects do not & \code{array}, NumPy arrays instead of lists of floats \\
Parallelism & the CPU applies one instruction to many values (SIMD), or many cores work at once & NumPy and pandas operations; multiprocessing \\
Concurrency & do other work while waiting for disk or network & threads, \code{asyncio} (Chapter 21) \\
\bottomrule
\end{tabularx}
\end{center}

% ------------------------------------------------------------
\section{Measure Timing}

\begin{center}
\begin{tabularx}{\linewidth}{@{}l X X@{}}
\toprule
\textbf{Tool} & \textbf{Good for} & \textbf{Note} \\
\midrule
\code{time.time()} & wall-clock epoch seconds & can jump if the system clock changes; the book uses it \\
\code{time.perf\_counter()} & timing a block of code & best clock for intervals \\
\code{timeit} & tiny snippets, many repeats & runs the code for you \\
\code{cProfile} & which function is slow & per-function table \\
\code{tracemalloc} & which code uses memory & peak and current bytes \\
\bottomrule
\end{tabularx}
\end{center}

\subsection{A stopwatch with perf\_counter}

\begin{pycode}
import time

t0 = time.perf_counter()
time.sleep(0.2)                          # the work being timed
elapsed = time.perf_counter() - t0
print(f"{elapsed:.2f} s")
\end{pycode}
\begin{rawcode}
0.20 s
\end{rawcode}

Run it several times: tiny jobs vary a lot (other programs, caches, power saving). One measurement
proves little.

\subsection{timeit}

\code{timeit.timeit(code, number=n)} runs the code \code{n} times and returns the \emph{total}
seconds for all runs (a \code{float}). The default \code{number} is one million, so set it
yourself for anything slow. \code{repeat} gives several totals; take the \textbf{minimum}, which has
the least noise. You can pass a string, or any callable.

\begin{pycode}
from timeit import timeit, repeat

total = timeit("sum(range(100))", number=1000)       # float, seconds
runs = repeat("sum(range(100))", number=1000, repeat=5)
best = min(runs) / 1000                             # seconds per call
print(type(total).__name__, len(runs))
print(f"best per call: {best * 1e6:.1f} us")        # varies per machine

data = list(range(1000))
t = timeit(lambda: sorted(data), number=100)        # pass a callable
print(f"{t:.4f} s for 100 sorts")
\end{pycode}
\begin{rawcode}
float 5
best per call: 0.4 us
0.0003 s for 100 sorts
\end{rawcode}
\noindent\textit{Your numbers will differ.}

\subsection{A decorator and a context manager}

For your own programs, wrap timing in a reusable tool. \code{functools.wraps} keeps the
function's name; \code{try/finally} prints even if the function raises.

\begin{pycode}
import time
from contextlib import contextmanager
from functools import wraps

def timed(func):
    @wraps(func)
    def inner(*args, **kwargs):
        t0 = time.perf_counter()
        try:
            return func(*args, **kwargs)
        finally:
            print(f"{func.__name__}: {time.perf_counter() - t0:.3f} s")
    return inner

@contextmanager
def stopwatch(label):
    t0 = time.perf_counter()
    yield
    print(f"{label}: {time.perf_counter() - t0:.3f} s")

@timed
def embed_batch(n):                 # stand-in for an embedding call
    time.sleep(0.05 * n)
    return [0.0] * n

vecs = embed_batch(2)
with stopwatch("parse"):
    time.sleep(0.1)
\end{pycode}
\begin{rawcode}
embed_batch: 0.100 s
parse: 0.100 s
\end{rawcode}

\begin{trap}
Do not time code with \code{print} in the middle, with the first run only (caches and imports make it
slower), or with a stdout redirect that does I/O. And never compare timings from two different
machines or from your laptop on battery. Compare two versions on the same machine in the same run.
\end{trap}

% ------------------------------------------------------------
\section{Profile}

Timers tell you \emph{how long}. A profiler tells you \emph{where}. \code{cProfile} (written in C,
low overhead) counts calls and times every function. Below, a pipeline hides a slow
de-duplication step; the profile exposes it. I ran this; file paths are shortened.

\begin{pycode}
import cProfile
import pstats

def tokenize(text):
    return text.split()

def slow_unique(words):
    seen = []
    for w in words:
        if w not in seen:           # linear search in a list
            seen.append(w)
    return seen

def pipeline():
    text = " ".join(f"w{i % 700}" for i in range(20000))
    return slow_unique(tokenize(text))

prof = cProfile.Profile()
prof.runcall(pipeline)
pstats.Stats(prof).sort_stats("cumulative").print_stats(5)
\end{pycode}
\begin{rawcode}
         20707 function calls in 0.037 seconds

   Ordered by: cumulative time

   ncalls  tottime  percall  cumtime  percall filename:lineno(function)
        1    0.000    0.000    0.037    0.037 pipeline.py:14(pipeline)
        1    0.033    0.033    0.033    0.033 pipeline.py:7(slow_unique)
        1    0.002    0.002    0.004    0.004 {method 'join' of 'str' objects}
    20001    0.003    0.000    0.003    0.000 pipeline.py:15(<genexpr>)
        1    0.000    0.000    0.000    0.000 pipeline.py:4(tokenize)
\end{rawcode}

\begin{center}
\begin{tabular}{@{}l l l@{}}
\toprule
\textbf{Column} & \textbf{Meaning} & \textbf{When to look} \\
\midrule
\code{ncalls} & number of calls & huge numbers: called too often \\
\code{tottime} & time inside the function itself & the real hotspot \\
\code{cumtime} & time including everything it calls & which top-level step is slow \\
\bottomrule
\end{tabular}
\end{center}

\begin{boss}[: measure before you optimize]
Everyone guesses. In the run above you might blame \code{tokenize} or the string building. The
profiler shows that one function, \code{slow\_unique}, owns about 90\% of the time.

%MD
%```mermaid
%flowchart LR
%    P["pipeline: 0.037 s total"] --> S["slow_unique: 0.033 s (89%)"]
%    P --> J["join + generator: 0.004 s (11%)"]
%    P --> T["tokenize: under 0.001 s"]
%```
%ENDMD
\begin{center}
\begin{tikzpicture}[font=\sffamily\small,>=Stealth]
  \node[w27n=cgray,minimum width=3.2cm] (p) at (0,0) {\code{pipeline}\\{\scriptsize 0.037 s total}};
  \node[w27n=cred,minimum width=3.2cm]  (s) at (5.2,1.1)  {\code{slow\_unique}\\{\scriptsize 0.033 s}};
  \node[w27n=cgold,minimum width=3.2cm] (j) at (5.2,-0.2) {\code{join} + generator\\{\scriptsize 0.004 s}};
  \node[w27n=cgreen,minimum width=3.2cm](t) at (5.2,-1.3) {\code{tokenize}\\{\scriptsize $<$0.001 s}};
  \draw[->,thick] (p) -- (s); \draw[->,thick] (p) -- (j); \draw[->,thick] (p) -- (t);
  % proportional bars
  \fill[cred!70]   (8.0,1.0) rectangle (8.0+0.033*100,1.3);
  \fill[cgold!70]  (8.0,-0.3) rectangle (8.0+0.004*100,0.0);
  \fill[cgreen!70] (8.0,-1.4) rectangle (8.0+0.0005*100,-1.1);
  \node[font=\scriptsize,anchor=west] at (11.5,1.15) {89\%};
  \node[font=\scriptsize,anchor=west] at (8.5,-0.15) {11\%};
\end{tikzpicture}
\end{center}

The fix is a data-structure change, not clever code: membership in a list scans every item, in a set
it is a hash lookup. One line, found only by measuring:

\begin{pycode}
def fast_unique(words):
    return list(dict.fromkeys(words))     # keeps first-seen order, O(n)
\end{pycode}
\end{boss}

For memory, \code{tracemalloc} reports the peak bytes allocated while your code ran (measured
here, one million numbers):

\begin{pycode}
import tracemalloc
from array import array

def peak_mb(build):
    tracemalloc.start()
    keep = build()
    _, peak = tracemalloc.get_traced_memory()
    tracemalloc.stop()
    return peak / 1_000_000

N = 1_000_000
print("list of ints :", round(peak_mb(lambda: list(range(N))), 1), "MB")
print("array('i')   :", round(peak_mb(lambda: array("i", range(N))), 1), "MB")
gen_mb = peak_mb(lambda: sum(x for x in range(N)))
print("sum(generator):", round(gen_mb, 3), "MB")
\end{pycode}
\begin{rawcode}
list of ints : 40.0 MB
array('i')   : 4.1 MB
sum(generator): 0.0 MB
\end{rawcode}

A generator never holds the whole sequence, which matters when you stream a huge file of
documents.

% ------------------------------------------------------------
\section{Algorithms and Data Structures}

``Better'' can mean faster, shorter, clearer or lighter on memory. Compare with data. All numbers
below are measured on this machine (Python 3.12); yours will differ, but the \emph{ratios} are the
lesson.

\begin{pycode}
from timeit import timeit

def with_append():
    out = []
    for v in range(1000):
        out.append(v * 2)
    return out

def with_comp():
    return [v * 2 for v in range(1000)]

print("loop :", round(timeit(with_append, number=2000), 3))
print("comp :", round(timeit(with_comp, number=2000), 3))

ids_list = list(range(5000))
ids_set = set(ids_list)
print("list in:", round(timeit(lambda: 4999 in ids_list, number=2000), 4))
print("set in :", round(timeit(lambda: 4999 in ids_set, number=2000), 6))

parts = [str(i) for i in range(2000)]
def plus():
    s = ""
    for p in parts:
        s += p
    return s
print("join :", round(timeit(lambda: "".join(parts), number=500), 4))
print("plus :", round(timeit(plus, number=500), 4))
\end{pycode}
\begin{rawcode}
loop : 0.052
comp : 0.043
list in: 0.0342
set in : 6e-05
join : 0.0024
plus : 0.0488
\end{rawcode}

%MD
%```text
%Time for the whole run (measured, log scale)
%append loop    | ###################################  0.052 s
%comprehension  | ##################################   0.043 s
%in list        | #################################    0.0342 s
%in set         | ########                             0.00006 s
%join           | ###########################          0.0024 s
%plus-equals    | ###################################  0.0488 s
%```
%ENDMD
\begin{center}
\begin{tikzpicture}
\begin{axis}[xbar,xmode=log,xmin=10,xmax=1e6,width=11cm,height=6.4cm,
  symbolic y coords={append loop,comprehension,in list,in set,join,plus-equals},
  ytick=data,y dir=reverse,xlabel={time for the whole run, microseconds (log scale, measured)},
  bar width=8pt,enlarge y limits=0.12,nodes near coords,point meta=explicit symbolic,
  every node near coord/.append style={font=\scriptsize,anchor=west}]
\addplot[fill=cblue!60,draw=cblue] coordinates {
  (52000,append loop) [0.052 s] (43000,comprehension) [0.043 s] (34200,in list) [0.0342 s]
  (60,in set) [0.00006 s] (2400,join) [0.0024 s] (48800,plus-equals) [0.0488 s]};
\end{axis}
\end{tikzpicture}
\end{center}

\begin{center}
\begin{tabularx}{\linewidth}{@{}l l X@{}}
\toprule
\textbf{Choice} & \textbf{Measured gap} & \textbf{Why} \\
\midrule
comprehension vs append loop & about 1.2$\times$ here & fewer bytecode steps; the book saw about 2$\times$ on its machine, so gaps change between Python versions \\
\code{in} on a set vs a list & about 500$\times$ & hash lookup vs scanning every item \\
\code{"".join(parts)} vs repeated \code{+=} & about 20$\times$ & join measures once and copies once \\
\bottomrule
\end{tabularx}
\end{center}

\begin{gamerule}
The big wins come from changing the \textbf{algorithm or data structure} (list to set, repeated work
to cached work), not from micro-tweaks. Choose the readable form first; optimize only the measured
hotspot.
\end{gamerule}

\subsection{Arrays versus lists and tuples}

Contiguous, same-type data is friendlier to the CPU cache. The standard \code{array} module
packs numbers tightly (measured above: 4.1 MB against 40.0 MB for a million ints). NumPy goes
further and runs the loop in compiled code (Chapter 25). NumPy is not installed here, so this is a
syntax-checked sketch:

\begin{pycode}
import numpy as np

values = list(range(1_000_000))
doubled_slow = [v * 2 for v in values]       # Python loop

arr = np.arange(1_000_000)
doubled_fast = arr * 2                       # compiled loop, no Python per item
\end{pycode}

\begin{typecheck}
The speed of NumPy and pandas comes from keeping work \emph{out of Python's per-item loop}. If you
call \code{.apply} or loop over \code{df.iterrows()}, you lose the benefit. Prefer whole-column
operations.
\end{typecheck}

% ------------------------------------------------------------
\section{Cache}

If the same call keeps producing the same answer, compute it once. \code{functools.cache} (unbounded)
and \code{lru\_cache(maxsize=N)} (keeps the N most recently used) store results keyed by the
arguments. This is ideal for expensive, repeatable work, such as embedding the same text twice.

%MD
%```mermaid
%flowchart LR
%    CALL["call f(args)"] --> Q{"in cache?"}
%    Q -->|yes| HIT["return saved result (fast)"]
%    Q -->|no| MISS["compute, save, return (slow, first time only)"]
%```
%ENDMD
\begin{center}
\begin{tikzpicture}[font=\sffamily\small,>=Stealth]
  \node[w27n=cblue]  (call) at (0,0) {call \code{f(args)}};
  \node[w27n=cgold,minimum width=2.4cm] (q) at (4.0,0) {in cache?};
  \node[w27n=cgreen] (hit) at (8.0,0.9)  {return saved result\\{\scriptsize fast}};
  \node[w27n=cred]   (miss) at (8.0,-0.9) {compute, \textbf{save}, return\\{\scriptsize slow, first time only}};
  \draw[->,thick] (call) -- (q);
  \draw[->,thick,cgreen] (q) -- node[above,font=\scriptsize]{yes} (hit);
  \draw[->,thick,cred] (q) -- node[below,font=\scriptsize]{no} (miss);
\end{tikzpicture}
\end{center}

\begin{pycode}
import time
from functools import lru_cache, cache

calls = 0

@cache
def embed(text):                 # pretend this calls an embedding API
    global calls
    calls += 1
    time.sleep(0.05)
    return len(text)

t0 = time.perf_counter()
for t in ["a", "bb", "a", "bb", "a"]:
    embed(t)
print(calls, f"{time.perf_counter() - t0:.2f}")
print(embed.cache_info())

@lru_cache(maxsize=2)            # small limit: evicts the oldest
def f(x): return x
for x in (1, 2, 3, 1):
    f(x)
print(f.cache_info())

try:
    @cache
    def g(xs): return len(xs)
    g([1, 2])
except TypeError as e:
    print("TypeError:", e)
\end{pycode}
\begin{rawcode}
2 0.10
CacheInfo(hits=3, misses=2, maxsize=None, currsize=2)
CacheInfo(hits=0, misses=4, maxsize=2, currsize=2)
TypeError: unhashable type: 'list'
\end{rawcode}

Five calls, only two real computations. In the second cache, the last call to \code{f(1)} missed
because \code{1} had already been evicted.

\begin{trap}
Arguments must be \textbf{hashable}: pass a tuple or string, not a list. Caching ties up memory
(use \code{lru\_cache(maxsize=...)} for unbounded inputs). Never cache a function with side effects
or one whose answer changes over time (the current time, a live API result).
\end{trap}

% ------------------------------------------------------------
\section{Concurrency for waiting}

Chapter 21's lesson applies directly to AI apps: most of your time is spent waiting for model and
network calls. Overlap the waits. Measured here with a fake call that sleeps 0.1 s, ten of them:

\begin{pycode}
import asyncio
import time
from concurrent.futures import ThreadPoolExecutor

def fake_call(i):
    time.sleep(0.1)                  # stands in for a network wait
    return i

async def fake_acall(i):
    await asyncio.sleep(0.1)
    return i

t0 = time.perf_counter()
seq = [fake_call(i) for i in range(10)]
print(f"sequential: {time.perf_counter() - t0:.1f} s")

t0 = time.perf_counter()
with ThreadPoolExecutor(max_workers=10) as pool:
    res = list(pool.map(fake_call, range(10)))
print(f"threads   : {time.perf_counter() - t0:.1f} s")

async def main():
    return await asyncio.gather(*(fake_acall(i) for i in range(10)))
t0 = time.perf_counter()
res2 = asyncio.run(main())
print(f"asyncio   : {time.perf_counter() - t0:.1f} s", res2 == seq)
\end{pycode}
\begin{rawcode}
sequential: 1.0 s
threads   : 0.1 s
asyncio   : 0.1 s True
\end{rawcode}

%MD
%```text
%Ten 0.1 s waits (measured)
%sequential | ################################  1.0 s
%threads    | ###                               0.1 s
%asyncio    | ###                               0.1 s
%```
%ENDMD
\begin{center}
\begin{tikzpicture}
\begin{axis}[ybar,symbolic x coords={sequential,threads,asyncio},xtick=data,ymin=0,ymax=1.25,
  ylabel={seconds},width=8.5cm,height=5cm,bar width=24pt,nodes near coords,enlarge x limits=0.3,
  every node near coord/.append style={font=\scriptsize,/pgf/number format/fixed,/pgf/number format/precision=1},
  title={\scriptsize ten 0.1 s waits (measured)},title style={font=\sffamily}]
\addplot[fill=cgreen!60,draw=cgreen] coordinates {(sequential,1.0) (threads,0.1) (asyncio,0.1)};
\end{axis}
\end{tikzpicture}
\end{center}

\begin{typecheck}
Threads and \code{asyncio} help with \emph{waiting} (I/O). They do not speed up pure number
crunching in CPython, because of the global interpreter lock. For CPU-bound work use NumPy, a
compiled tool below, or multiple processes.
\end{typecheck}

% ------------------------------------------------------------
\section{Going faster than CPython}

When the algorithm is right and the hotspot is still CPU-bound Python, you can compile it. First,
how Python normally runs, and where each tool plugs in:

%MD
%```mermaid
%flowchart LR
%    subgraph CPython
%        A1["source"] --> A2["bytecode"] --> A3["interpreter loop: steps through bytecode"]
%    end
%    subgraph JIT["JIT (PyPy, Numba, 3.13 JIT)"]
%        B1["source"] --> B2["bytecode"] --> B3["machine code, compiled while running"]
%    end
%    subgraph AOT["Ahead of time (Cython, C/Rust extensions, Mojo)"]
%        C1["typed source"] --> C2["compiler"] --> C3["machine code: import and call"]
%    end
%```
%ENDMD
\begin{center}
\begin{tikzpicture}[font=\sffamily\small,>=Stealth]
  % row 1: CPython
  \node[anchor=east,font=\sffamily\bfseries\scriptsize] at (-0.2,2.2) {CPython};
  \node[w27n=cgray]   (a1) at (1.2,2.2) {source};
  \node[w27n=cblue]   (a2) at (4.2,2.2) {bytecode};
  \node[w27n=cgold,minimum width=3.4cm] (a3) at (8.4,2.2) {interpreter loop\\{\scriptsize steps through bytecode}};
  \draw[->,thick] (a1) -- (a2); \draw[->,thick] (a2) -- (a3);
  % row 2: JIT
  \node[anchor=east,font=\sffamily\bfseries\scriptsize] at (-0.2,0.8) {JIT};
  \node[w27n=cgray]   (b1) at (1.2,0.8) {source};
  \node[w27n=cblue]   (b2) at (4.2,0.8) {bytecode};
  \node[w27n=cgreen,minimum width=3.4cm] (b3) at (8.4,0.8) {machine code\\{\scriptsize compiled while running}};
  \draw[->,thick] (b1) -- (b2); \draw[->,thick] (b2) -- (b3);
  \node[font=\scriptsize,text=cgray,anchor=west] at (10.3,0.8) {PyPy, Numba, 3.13 JIT};
  % row 3: AOT
  \node[anchor=east,font=\sffamily\bfseries\scriptsize] at (-0.2,-0.6) {Ahead of time};
  \node[w27n=cgray]   (c1) at (1.2,-0.6) {typed source\\{\scriptsize .pyx, Rust, Mojo}};
  \node[w27n=cpurple] (c2) at (4.6,-0.6) {compiler};
  \node[w27n=cgreen,minimum width=3.4cm] (c3) at (8.4,-0.6) {machine code\\{\scriptsize import and call}};
  \draw[->,thick] (c1) -- (c2); \draw[->,thick] (c2) -- (c3);
  \node[font=\scriptsize,text=cgray,anchor=west] at (10.3,-0.6) {Cython, C/Rust ext.};
\end{tikzpicture}
\end{center}

You can see CPython's bytecode with the \code{dis} module (output from Python 3.12):

\begin{pycode}
import dis
dis.dis("a + b")
\end{pycode}
\begin{rawcode}
  0           0 RESUME                   0

  1           2 LOAD_NAME                0 (a)
              4 LOAD_NAME                1 (b)
              6 BINARY_OP                0 (+)
             10 RETURN_VALUE
\end{rawcode}

\begin{center}
\begin{tabularx}{\linewidth}{@{}l X X X@{}}
\toprule
\textbf{Tool} & \textbf{How it works} & \textbf{You change} & \textbf{Good for} \\
\midrule
Cython & Python with optional C types, compiled to C & add \code{cdef} types, build step & numeric loops (large speedups when typed) \\
C or Rust extension & write the hot function in C/Rust, call from Python & write in another language & libraries; Rust is memory-safe. Pydantic's core and \code{uv} are Rust \\
Taichi & JIT-compiles Python-like kernels & decorate kernels & numeric and graphics code \\
PyPy & alternative interpreter with a JIT & run your program with \code{pypy} & long-running pure-Python code; check C-extension support \\
Numba & JIT for numeric functions & add \code{@njit} & loops over numbers and NumPy arrays \\
CPython 3.13 JIT & experimental built-in JIT & rebuild CPython with the option & experiments only \\
Mojo & new Python-like language with types, compiled & rewrite hot code in Mojo & ML kernels; still young \\
\bottomrule
\end{tabularx}
\end{center}

\subsection{Cython}

Cython translates Python plus type declarations to C. Not run here (needs a C compiler and a build
step); illustrative \code{.pyx} file:

\begin{rawcode}
# sumsq.pyx
def sum_sq(int n):
    cdef long total = 0
    cdef int i
    for i in range(n):
        total += i * i
    return total
\end{rawcode}

\subsection{NumPy and SciPy}

The cheapest ``compiler'' is a library that already has compiled loops. NumPy and SciPy do the heavy
numeric work in C and Fortran (see Chapter 25), which is a big reason Python is used for data
science.

\subsection{C or Rust extensions}

Much of Python itself and its standard library is C wrapped for convenience. You can write your own
extension for a hotspot. A recent trend is Rust instead of C because it is fast and safer with
memory. Two well-known Python tools built this way: Pydantic's validation core, and the package
manager \code{uv}.

\subsection{Taichi, PyPy and Numba}

\textbf{CPython} is the standard interpreter (written in C; do not confuse it with Cython). It
translates your code to bytecode and steps through it. \textbf{PyPy} is another interpreter that
adds just-in-time compilation and is often much faster on CPU-bound pure Python. \textbf{Numba}
compiles chosen functions with a decorator. The first call includes compile time, so time later
calls:

\begin{pycode}
from numba import njit          # pip install numba

@njit                           # = jit(nopython=True)
def sum_sq(n):
    total = 0
    for i in range(n):
        total += i * i
    return total

sum_sq(10)                      # first call compiles; ignore its time
result = sum_sq(10_000_000)     # later calls run machine code
\end{pycode}

Syntax-checked only. A JIT helps with long loops over numbers. It will not help a tiny function
called once, or code that is mostly I/O or dictionaries of strings (the book's own tiny
\code{hypot} example showed little gain).

\subsection{The standard-library JIT}

Python 3.13 added an \emph{experimental} JIT, available only if you rebuild CPython with the right
flag. That version also has an experimental free-threaded build (no GIL). Neither is something to
rely on yet.

\subsection{Mojo}

Mojo is a newer language from Modular (announced in 2023), designed as a faster ``Python++'' for AI
work. Key points from the book:

\begin{itemize}[leftmargin=*]
\item \textbf{Why:} modern AI stacks are a sandwich: Python on top, C/C++/CUDA in the middle, new
hardware underneath. Mojo aims to cover all layers in one language.
\item \textbf{Who:} Chris Lattner, who earlier created LLVM, Clang, Swift and MLIR.
\item \textbf{Design:} compiles to machine code; type hints are mandatory (CPython ignores them);
\code{struct} groups data; \code{fn} replaces \code{def} for strictly typed functions; it can run
existing Python.
\item \textbf{Claims:} one demo was 35{,}000$\times$ faster than naive Python; another about 60$\times$
and twice NumPy. Treat such vendor and demo figures with caution and measure your own workload.
\item \textbf{Limits:} young; Python compatibility is still incomplete.
\end{itemize}

\begin{gamerule}
Climb the ladder in order and stop when it is fast enough: (1) measure, (2) better algorithm or data
structure, (3) cache or avoid repeated work, (4) use built-in or C-backed libraries (\code{sum},
\code{sorted}, NumPy, pandas), (5) overlap waiting with threads or \code{asyncio}, (6) compile the
hotspot (Numba, Cython, Rust), (7) different interpreter or language. Each step costs more effort
and flexibility than the previous one.
\end{gamerule}

% ------------------------------------------------------------
\section{Practice and wrap-up}

Performance work is detective work: measure, find the cause, change one thing, measure again.

\begin{practice}
\begin{enumerate}[leftmargin=*]
\item Why is \code{time.perf\_counter()} a better stopwatch than \code{time.time()}? What does
\code{timeit.repeat} return, and which value of it should you report?
\item Profile a function of your own. Which column tells you the hotspot, \code{tottime} or
\code{cumtime}? Why?
\item Change a list used for membership tests to a set and measure. Predict the ratio first.
\item Why does \code{@cache} fail for a function that takes a list? How would you fix the call?
\item Embedding 200 documents takes 0.1 s of waiting per call. Roughly how long does a sequential
loop take, and a pool of 20 threads? What changes if the work were CPU-bound instead?
\item Compare three ways to compute the sum of the first $n$ squares: a \code{for} loop, a generator
passed to \code{sum}, and the closed formula $n(n+1)(2n+1)/6$. Time them for $n=10^6$ and explain the
result.
\end{enumerate}
\end{practice}

\begin{achievement}{Performance detective}
You can now:
\begin{itemize}[leftmargin=*]
\item time code correctly with \code{perf\_counter} and \code{timeit}, and report a stable number;
\item read a \code{cProfile} table to find the real hotspot, and check memory with \code{tracemalloc};
\item replace a slow data structure or algorithm (list to set, \code{+=} to \code{join}, loop to comprehension);
\item add caching and concurrency, and know when each one helps;
\item explain where Numba, Cython, PyPy, Rust extensions and Mojo fit, and when to leave them alone.
\end{itemize}
\end{achievement}
